\documentclass[11pt,a4paper]{article}
\usepackage[utf8]{inputenc}
\usepackage{geometry}
\geometry{margin=1in}
\usepackage{lineno}
\usepackage{xcolor}
\usepackage{titlesec}
\usepackage{array}

% Command for redacting personal information with a black box
\newcommand{\redacted}{\rule{2.5cm}{2ex}}

\title{\textbf{Social Science Fair Interview 01 \\ Transcript: Value and Pressure of Football Defenders}}
\author{\textbf{Research Team Members:} Luke Sun, Kiara Wei, Winston Huang}
\date{May 28, 2026}

\begin{document}

\maketitle

\section*{Transcript Metadata}
\begin{tabular}{ll}
    \textbf{Project Theme:} & Discussion on the Value and Pressure of Football Defenders \\
    \textbf{Date of Interview:} & Thursday, May 28, 2026 \\
    \textbf{Time of Interview:} & 19:52 -- 20:04 (GMT+08) \\
    \textbf{Interviewer Identifier:} & Speaker 100 Luke Sun\\
    \textbf{Interviewee Identifier:} & Speaker 200 (\redacted) \\
    \textbf{Interviewee Subject Profile:} & School Team Athlete (Right-Back)
\end{tabular}

\vspace{2em}
\hrule
\vspace{1.5em}

\section*{Official Verbatim Transcript}

\begin{linenumbers}

\noindent \textbf{Speaker 100 [00:01]:} Alright, you are \redacted, correct? 

\vspace{0.5em}
\noindent \textbf{Speaker 200 [00:04]:} Yes, I am. 

\vspace{0.5em}
\noindent \textbf{Speaker 100 [00:06]:} Okay. What position do you currently play on the school team? 

\vspace{0.5em}
\noindent \textbf{Speaker 200 [00:10]:} I mainly play right-back. That position, yes. 

\vspace{0.5em}
\noindent \textbf{Speaker 100 [00:14]:} And how do you define your core responsibility on the pitch? What do you primarily do, roughly? 

\vspace{0.5em}
\noindent \textbf{Speaker 200 [00:21]:} Simply put, it's getting the ball, kicking it away, and passing it to the players upfront so they can attack. To put it more complexly, when an opportunity arises, I intercept the ball and organize the transition from the backfield to the frontfield. Additionally, when the opponent attacks, we form a second line of defense to disrupt their offense. 

\vspace{0.5em}
\noindent \textbf{Speaker 100 [00:45]:} In the last match, how much of a decisive role do you think you played in the team's victory or defeat? 

\vspace{0.5em}
\noindent \textbf{Speaker 200 [00:52]:} I feel my role in the outcome wasn't huge, but I still feel there was some. Well, that's a bit of nonsense first, but I think the main reason is that I played for a relatively short time; I played for a bit and then got substituted off. However, I think I was useful because when I was playing, I helped absorb a lot of pressure for the center-backs, which maintained our defensive advantage. 

\vspace{0.5em}
\noindent \textbf{Speaker 100 [01:28]:} Understood. When your mistake leads to the team conceding a goal, do you feel intense pressure from your teammates, the coach, or the spectators? As a defender, because you missed a ball and caused a concession, you must have experienced that, right? 

\vspace{0.5em}
\noindent \textbf{Speaker 200 [01:45]:} Yes. 

\vspace{0.5em}
\noindent \textbf{Speaker 100 [01:46]:} Do you feel the pressure from them is heavy? 

\vspace{0.5em}
\noindent \textbf{Speaker 200 [01:51]:} I feel the pressure from them is quite heavy. Although I am inherently a pressure-free person—when others put pressure on me, I don't really convert it into internal friction or things like that—I feel their pressure helps me realize where I need to improve, rather than making me anxious about which direction I haven't done better in. 

\vspace{0.5em}
\noindent \textbf{Speaker 100 [02:20]:} Okay. When the team wins, to what extent do you feel the credit is attributed to your performance on the defensive line? And to what extent is it seen as the performance of the attacking line, or the result of a collective team effort? 

\vspace{0.5em}
\noindent \textbf{Speaker 200 [02:33]:} I feel it's more like the result of a collective team effort. Because in a football match, both defense and offense are very important, right? Although some Champions League teams might not agree, I still believe that winning is the collective outcome of the entire team acting as a unified entity, including the substitutes on the bench. Their encouragement gives the players on the field more fighting spirit, which also plays a certain role. So, no matter who it is, everyone is indispensable. 

\vspace{0.5em}
\noindent \textbf{Speaker 100 [03:10]:} Um, okay. Do you feel that defensive players face harsher criticism and negative attention than attacking players? 

\vspace{0.5em}
\noindent \textbf{Speaker 200 [03:21]:} I feel this is a very common issue. For example, if you, as the defense, lose the ball, the opponent has a very high probability of scoring. But if you, as the offense, lose the ball, it merely gives the opponent an attacking opportunity. It doesn't directly lead to a highly irreversible consequence for the team. To give a more extreme example, when the game is intense and tied, going into extra time, if the backfield loses a ball at this moment, leading to an opponent's goal, it is very difficult for our side to score again in such a short time. But if the offense loses the ball, the defense can still win it back. 

\vspace{0.5em}
\noindent \textbf{Speaker 100 [04:13]:} In post-match summaries or daily discussions, do you talk more often about statistics like goals and assists, or data like passing under pressure or critical tackles from defenders? 

\vspace{0.5em}
\noindent \textbf{Speaker 200 [04:26]:} I feel if a critical tackle looks exceptionally cool, people will definitely talk about it, but generally speaking, it's still about things like goals. 

\vspace{0.5em}
\noindent \textbf{Speaker 100 [04:35]:} Um, do you feel your defensive contribution is somewhat undervalued because it's hard to quantify like the number of goals? 

\vspace{0.5em}
\noindent \textbf{Speaker 200 [04:45]:} Not really, because I feel if I defend and our team wins, I feel my defense at least helped my team. If I keep defending and the opponent pours in 5 goals, then I feel I'm just purely bad, and there's nothing more to say. 

\vspace{0.5em}
\noindent \textbf{Speaker 100 [05:03]:} In the market, they usually count goals and assists. You rarely see people counting clearances—there are some, but the attention isn't that high—and things like tackle counts are even rarer. Do you think this current situation leads to the undervaluation of players' worth? That is, will players who lean towards being defensive-minded be undervalued? 

\vspace{0.5em}
\noindent \textbf{Speaker 200 [05:26]:} I feel they will be. For instance—actually, I don't have that many examples because I watch less football now—but I remember, for example, Piqué from Barcelona in 16/17, he was a center-back. But he defended exceptionally well as a center-back. Although it was largely attributed to his tall stature and excellent technique, I personally feel their success in that season drew heavily from the success of the backline. However, over those years, the main credit went to Messi. It's not that Messi wasn't good; Messi indeed could single-handedly carry the attacking side and score many goals, which is truly strong. But you have to know that in football, the pressure exerted on the defense is usually greater than the pressure applied to the offense. Because when we play, the offense losing two or three chances to get one or two goals is very worthwhile. But for the defense, if the opponent gets two or three chances and scores a goal, it feels terrible. 

\vspace{0.5em}
\noindent \textbf{Speaker 100 [06:38]:} Oh, right. Then do you think this potential undervaluation, whether amateur or professional, affects a player's status or confidence in the team? For example, you potentially face this kind of undervaluation, right? Do you feel that way? 

\vspace{0.5em}
\noindent \textbf{Speaker 200 [06:58]:} Actually, I feel my situation is quite special, but I do feel a little bit, because I transitioned from midfield. 

\vspace{0.5em}
\noindent \textbf{Speaker 100 [07:04]:} Then do you think this potential undervaluation has affected your status or confidence in the team? 

\vspace{0.5em}
\noindent \textbf{Speaker 200 [07:13]:} It hasn't really affected my confidence. As for status, to put it bluntly, I didn't have much status to begin with. But personally, I think the change in status wasn't very big. Of course, I feel it decreased slightly because, in terms of impact on the game, I personally believe transitioning from midfield to right-back reduced my ability to influence the pitch. Although this change might be better for the team, so it doesn't matter. 

\vspace{0.5em}
\noindent \textbf{Speaker 100 [07:48]:} Um, okay, okay. As mentioned earlier, a backfield player losing the ball often leads to a direct goal, while a frontfield player losing the ball can be won back, representing a relatively lower risk. We call this the asymmetry of risk. Do you think there is any possibility to change this reality through specific measures, plans, or policies? For example, promoting the art of defending by defenders more on the internet, or teaching it more in offline football classrooms, encouraging less individual heroism and more teamwork, emphasizing that one cannot just attack but must also balance defense. Do you think there are feasible measures or plans like this? 

\vspace{0.5em}
\noindent \textbf{Speaker 200 [08:42]:} I feel if we want to balance the external pressure received by both sides, it might be possible. For example, if one day, as you mentioned, we promote from the opposite direction that the consequences of a forward losing the ball are more severe, and conduct campaigns through public opinion, then the pressure itself will become more balanced. But if you are talking about reducing the pressure on defenders, I feel the change won't be significant. Because a defender losing the ball and causing a defeat is a hard event. You basically cannot change this fact through external promotion. 

\vspace{0.5em}
\noindent \textbf{Speaker 200 [09:27]:} Right? Because it is inherently the fact that a goal was scored. Structurally speaking, defending is easier than attacking. So when you are attacking, losing the ball a bit more is considered normal by people. If you try to level it out by adjusting external expectations, to be honest, the gap between the blame carried for conceding a goal and the blame carried for losing the ball in the frontfield will still remain. 

\vspace{0.5em}
\noindent \textbf{Speaker 100 [10:03]:} Then are there any specific measures to level out this expectation? Have you thought of any? 

\vspace{0.5em}
\noindent \textbf{Speaker 200 [10:12]:} The first one that comes to mind right now is using the previous method of online public opinion. Another way, I feel, is to improve the overall level of the defensive line itself. For example, a three-man backline and a four-man backline are different. From another perspective, if a three-man backline concedes a goal, the blame is shared equally among these three people—excluding individual situations of who missed the ball, of course. If it's a four-man backline and a goal is conceded, this situation will be attributed to either the opponent being too strong, in which case it doesn't matter if there are three or four people, or all four of us sharing the blame simultaneously. In this case, the pressure will be distributed more evenly and won't be as high. 

\vspace{0.5em}
\noindent \textbf{Speaker 100 [11:09]:} So what are the specific measures? 

\vspace{0.5em}
\noindent \textbf{Speaker 200 [11:11]:} First, have defenders train more. Second, add more defenders. 

\vspace{0.5em}
\noindent \textbf{Speaker 100 [11:18]:} Add more defenders. 

\vspace{0.5em}
\noindent \textbf{Speaker 200 [11:20]:} That means strengthening the defensive capability of the team itself. By changing the hard strength, we reduce the probability of being overwhelmed by pressure. 

\vspace{0.5em}
\noindent \textbf{Speaker 100 [11:33]:} Okay, okay, okay, thank you. That's all.

\end{linenumbers}

\vspace{1em}

\textit{Translated from Chinese.}

\end{document}